%% Dans ce fichier, les �ventuels accents sont cod�s
%% en ISO-8859-15 (aka latin9)
%% La ligne suivante donne trois exemples 
%% e accent aigu : �, e accent grave : �, e-dans-l'o : �
%%%%%ENTREES%%%%
\begin{tabarg}
\multicolumn{3}{l}{\titretab {Entr\'ees}}\\
\hline
$P_s$ & Vecteur, \texttt{float} &  Pression statique de l'air ($hPa$)\\
$T_s$ & Vecteur, \texttt{float} & Temp�rature statique de l'air ($K$ ou $^\circ C$)\\
$V_p$ & Vecteur, \texttt{float} & Vitesse propre ($m/s$)\\
$P_w$ & Vecteur, \texttt{float} & Puissance fournie \`a la sonde ($W$)\\
$T_w$ & Vecteur, \texttt{float} & Temp\'erature de la sonde ($K$ ou $^\circ C$)\\
$K_1$ & Coeff., \texttt{float} & Coefficient de calcul de la puissance s\`eche Robust\\
$K_2$ & Coeff., \texttt{float} & Coefficient de calcul de la puissance s\`eche Robust\\
\end{tabarg}

\medskip
%%%%%SORTIES%%%%
\begin{tabarg}
\multicolumn{3}{l}{\titretab {Sorties}}\\
\hline
$LWC$ & Vecteur, \texttt{float} & Contenu en eau liquide ($g/m^3$)\\
\end{tabarg}
